This section records the \emph{flow-time quasi-locality and derivative bounds}
for the flowed composites $\mathcal O_{S,t}$ and $\mathcal O_{Q,t}$, in a form uniform in $a\to 0$ on the AF window.
The deterministic hypotheses used here (a local chart/boundedness regime for the DeTurck flow) are verified on the
good-flow chart event in Section~\ref{sec:good_flow_chart} (Lemma~\ref{lem:chart_persist}), and the contribution of the
complement is absorbed in expectations by truncation stability (Corollary~\ref{cor:superpoly_flow}).

\paragraph{Important correction (hardening).}
The gradient flow is \emph{parabolic}; it does \emph{not} have finite propagation speed.
Thus ``strict support in $B(x,c\sqrt t)$'' is false. What is true (and sufficient for the dyadic polymer/RG framework) is:
\begin{itemize}
\item \textbf{quasi-locality}: influence of initial data at distance $r$ is suppressed like $\exp(-c r^2/t)$;
\item \textbf{derivative bounds}: functional derivatives of the flow map and of $\mathcal O_{S,t},\mathcal O_{Q,t}$ satisfy Gaussian off-diagonal bounds with
explicit $t$-scaling factors;
\item \textbf{local truncation}: by cutting off outside a ball of radius $R\sqrt t$ one incurs an error $\lesssim e^{-cR^2}$, which can be made $\le t^\eta$ by choosing $R\sim \sqrt{\log(1/t)}$.
\end{itemize}

This module replaces 's strict-support formulation with the correct quasi-local statement and makes the constants explicit.

We work in Euclidean $\mathbb R^4$ (or a periodic box of side $L$ with $t\le t_\ast\ll L^2$ so boundary effects are negligible).
Let $A_\mu(x)\in\mathfrak g$ be the initial gauge potential and $B_\mu(t,x)$ the gradient-flowed field for $t>0$.

The standard (continuum) Yang--Mills gradient flow is
\begin{equation}\label{eq:cov_flow}
\partial_t B_\mu \ =\ D_\nu G_{\nu\mu},\qquad B_\mu(0,\cdot)=A_\mu,
\end{equation}
where $D_\nu=\partial_\nu+[B_\nu,\cdot]$ and $G_{\mu\nu}=\partial_\mu B_\nu-\partial_\nu B_\mu+[B_\mu,B_\nu]$.

To access parabolic estimates we use the DeTurck gauge-fixing term:
\begin{equation}\label{eq:deturck_flow}
\partial_t B_\mu \ =\ D_\nu G_{\nu\mu} + D_\mu(\partial_\nu B_\nu),\qquad B_\mu(0,\cdot)=A_\mu.
\end{equation}

\begin{proposition}[Gauge equivalence and invariance of gauge-invariant composites]\label{prop:deturck_equiv}
Solutions of \eqref{eq:cov_flow} and \eqref{eq:deturck_flow} are related by a (time-dependent) gauge transformation $g(t,x)\in G$:
\[
B^{\mathrm{cov}}_\mu(t)=g(t)\,B^{\mathrm{det}}_\mu(t)\,g(t)^{-1}-\partial_\mu g(t)\,g(t)^{-1}.
\]
Consequently, gauge-invariant local observables built from $G_{\mu\nu}(t)$, in particular
$\mathcal O_{S,t}(x)=\tfrac14\mathrm{tr}\,G_{\mu\nu}(t,x)G_{\mu\nu}(t,x)$ and
$\mathcal O_{Q,t}(x)=\tfrac14\mathrm{tr}\,G_{\mu\nu}(t,x)\widetilde G_{\mu\nu}(t,x)$,
have identical values under the two flows.
\end{proposition}

\paragraph{Comment.}
This is the standard DeTurck trick for Ricci/YM flows: the gauge-fixing term generates a compensating gauge transformation.
Thus we may work with \eqref{eq:deturck_flow} without loss for gauge-invariant observables.

To obtain uniform constants we work on a chart-controlled regime (consistent with the patched RG philosophy of Section~\ref{sec:patching}).

\begin{proposition}[Bounded-regime hypothesis (verified on the good-flow chart event)]\label{ass:bounded_regime}
Fix $t_\ast>0$. On the good-flow chart event $G_{\Lambda}(\varepsilon)$ of Definition~\ref{def:good_plaq} with initial chart extraction as in
Lemma~\ref{lem:tree_chart_local}, the DeTurck-flow solution satisfies the scale-consistent bound
\begin{equation}\label{eq:bounds_B}
\sup_{0<t\le t_\ast}\Big(\|B(t)\|_{L^\infty}+ t^{1/2}\|\nabla B(t)\|_{L^\infty}\Big)\ \le\ M_0+M_1.
\end{equation}
In lattice language, this corresponds to a fixed exponential chart bound on link angles and plaquette angles (e.g.\ $\le\delta<\pi/4$) on the relevant scales.
\end{proposition}

\paragraph{Where this comes from.}
Lemma~\ref{lem:chart_persist} (Section~\ref{sec:good_flow_chart}) gives a deterministic sufficient condition for
\eqref{eq:bounds_B}: on $G_{\Lambda}(\varepsilon)$ the flow stays inside a fixed exponential chart up to time $t_\ast$ and the parabolic estimates
control $\|B(t)\|_\infty$ and $t^{1/2}\|\nabla B(t)\|_\infty$ uniformly on $(0,t_\ast]$.
In the main proof we invoke reflection positivity/OS only on unflowed positive-time observables; when flowed insertions appear in expectations,
we enforce $G_{\Lambda}(\varepsilon)$ locally and absorb the complement by Corollary~\ref{cor:superpoly_flow}.

Let $\mathcal F_t$ denote the DeTurck flow map $A\mapsto B(t)$.
Let $\psi_\mu(t,x)$ denote the first variation of $B_\mu(t,x)$ under an infinitesimal perturbation of the initial data $A$.

\begin{lemma}[Linearized DeTurck flow is uniformly parabolic]\label{lem:lin_parabolic}
Under Proposition~\ref{ass:bounded_regime}, the first variation $\psi$ satisfies a linear parabolic system of the form
\begin{equation}\label{eq:lin_system}
\partial_t \psi \ =\ \Delta \psi + \sum_{\alpha} a_\alpha(t,x)\,\partial_\alpha \psi + b(t,x)\psi,
\end{equation}
where $\Delta$ is the Euclidean Laplacian acting componentwise and the coefficients satisfy the uniform bounds
\[
\begin{aligned}
\|a_\alpha(t,\cdot)\|_{L^\infty(\mathbb R^4)}&\le C_{a,0}(M_0),\\
\|b(t,\cdot)\|_{L^\infty(\mathbb R^4)}&\le C_{b,0}(M_0)+C_{b,1}(M_0,M_1)\,t^{-1/2},
\qquad 0<t\le t_\ast.
\end{aligned}
\]
\end{lemma}

\paragraph{Meaning.}
Equation \eqref{eq:lin_system} is uniformly parabolic; on any time slab $[s,t]\subset(0,t_\ast]$ the coefficients are bounded.
The possible $t^{-1/2}$ singularity at $t\downarrow 0$ is time-integrable and is treated by the parametrix/Volterra majorant
in Appendix~\ref{app:kernel_parametrix} (and uniformly in $a$ in Appendix~\ref{app:discrete_parametrix}).

\begin{proposition}[Gaussian kernel bounds for the linearized operator (proved in Appendix~\ref{app:kernel_parametrix} and Appendix~\ref{app:discrete_parametrix})]\label{ass:gaussian_kernel}
Let $K(t,s;x,y)$ be the fundamental solution (matrix kernel) of \eqref{eq:lin_system} on $0\le s<t\le t_\ast$.
Assume there exist constants $c_0,c_1>0$ and for each integer $m\ge 0$ constants $C_m$ (depending only on $d=4$ and the coefficient bounds
$C_{a,0},C_{b,0},C_{b,1}$ (and the corresponding finite set of derivative bounds used in the parametrix)) such that:
\begin{align}
\|K(t,s;x,y)\|_{\mathrm{op}}
&\le c_0\,(t-s)^{-2}\exp\!\Big(-c_1\,\frac{|x-y|^2}{t-s}\Big),\label{eq:K_gauss}\\
\|\nabla_x^m K(t,s;x,y)\|_{\mathrm{op}}
&\le C_m\,(t-s)^{-(2+m/2)}\exp\!\Big(-c_1\,\frac{|x-y|^2}{t-s}\Big).\label{eq:K_der_gauss}
\end{align}
\end{proposition}

On the lattice, an analogous statement holds for the discrete operator under uniform chart bounds; constants can be chosen independent of $a$ for $t\ge c a^2$.

Let $\delta A$ be a perturbation supported near a point $y$.

\begin{theorem}[First-variation quasi-locality of the flow map]\label{thm:first_var}
Under the bounded-regime hypothesis~\ref{ass:bounded_regime} and the Gaussian kernel bounds~\ref{ass:gaussian_kernel},
for $t\in(0,t_\ast]$ the Fr\'echet derivative of the flow map satisfies the kernel bound:
\[
\left\|\frac{\delta B_\mu(t,x)}{\delta A_\nu(y)}\right\|
\ \le\ c_0\,t^{-2}\exp\!\Big(-c_1\,\frac{|x-y|^2}{t}\Big),
\]
and for spatial derivatives:
\[
\left\|\nabla_x^m\frac{\delta B_\mu(t,x)}{\delta A_\nu(y)}\right\|
\ \le\ C_m\,t^{-(2+m/2)}\exp\!\Big(-c_1\,\frac{|x-y|^2}{t}\Big).
\]
Constants are those of Proposition~\ref{ass:gaussian_kernel} with $s=0$.
\end{theorem}

\begin{corollary}[First-variation bounds for field strength]\label{cor:first_var_G}
Under the same assumptions, the field strength variation satisfies for $m\ge 0$:
\[
\left\|\nabla_x^m\frac{\delta G_{\rho\sigma}(t,x)}{\delta A_\nu(y)}\right\|
\ \le\ C'_m\,t^{-(5/2+m/2)}\exp\!\Big(-c_1\,\frac{|x-y|^2}{t}\Big),
\]
where $C'_m$ depends on $C_{m+1}$ and $M_0$ (through commutator terms in $G$).
\end{corollary}

The polymer seminorms used in dyadic RG control derivatives up to a fixed finite order $r_{\max}$.
We therefore state a finite-order induction hypothesis.

\begin{proposition}[Finite-order higher derivative bounds (proved in Appendix~\ref{app:kernel_parametrix})]\label{ass:higher_der}
Fix $r_{\max}\ge 1$.
Assume that for each $1\le r\le r_{\max}$ there exist constants $C^{(B)}_{r,m}$ such that the $r$-th Fr\'echet derivative of the flow map obeys:
\begin{equation}\label{eq:high_der_B}
\left\|\nabla_x^m\frac{\delta^r B(t,x)}{\delta A(y_1)\cdots\delta A(y_r)}\right\|
\ \le\ C^{(B)}_{r,m}\, t^{-(2+m/2+r)}\,
\exp\!\Big(-c_1\,\frac{\mathrm{dist}(x;\{y_i\})^2}{t}\Big),
\end{equation}
where $\mathrm{dist}(x;\{y_i\})=\min_i |x-y_i|$.
\end{proposition}

\paragraph{Comment.}
The exponent $2+m/2+r$ is a safe (nonoptimal) power count; any explicit polynomial bound in $t^{-1/2}$ is sufficient for the A1 insertion norms,
because $r_{\max}$ is fixed and the ultimate SFTE remainder is controlled by extracting relevant/marginal parts and contracting irrelevants.
Proposition~\ref{ass:higher_der} is proved by differentiating the flow equation $r$ times and applying Duhamel/Gr\"onwall with the Gaussian kernel bounds.

\begin{corollary}[Higher derivative bounds for $\mathcal O_{S,t}$ and $\mathcal O_{Q,t}$]\label{cor:high_der_OSOQ}
Assume~\ref{ass:bounded_regime},~\ref{ass:gaussian_kernel}, and~\ref{ass:higher_der}.
Then for each $r\le r_{\max}$ there exist constants $C^{(S)}_{r}$ and $C^{(Q)}_{r}$ such that
\begin{align}
\left\|\frac{\delta^r \mathcal O_{S,t}(x)}{\delta A(y_1)\cdots\delta A(y_r)}\right\|
&\le C^{(S)}_{r}\, t^{-(5+r)}\,
\exp\!\Big(-c_1\,\frac{\mathrm{dist}(x;\{y_i\})^2}{t}\Big),\label{eq:der_OS}\\
\left\|\frac{\delta^r \mathcal O_{Q,t}(x)}{\delta A(y_1)\cdots\delta A(y_r)}\right\|
&\le C^{(Q)}_{r}\, t^{-(5+r)}\,
\exp\!\Big(-c_1\,\frac{\mathrm{dist}(x;\{y_i\})^2}{t}\Big).\label{eq:der_OQ}
\end{align}
The constants depend on $M_0,M_1$ through the coefficient bounds and on the fixed derivative order $r$.
\end{corollary}

\paragraph{Remark.}
The power $5+r$ is again safe (nonoptimal). The key structural feature is the Gaussian suppression $\exp(-c\,\mathrm{dist}^2/t)$.

We now convert quasi-locality into a localization statement usable in the dyadic polymer framework.

\begin{definition}[Ball truncation radius]\label{def:Rt}
For $R\ge 1$ define the ball $B_R(x,t):=\{y:|y-x|\le R\sqrt t\}$.
For a perturbation $\delta A$ supported in $B_R(x,t)^c$ define the Lipschitz seminorm
\[
\mathrm{Lip}^{(R)}_x(\mathcal O_{*,t})
:=
\sup_{\substack{\delta A\neq 0\\ \mathrm{supp}(\delta A)\subset B_R(x,t)^c}}
\frac{|\mathcal O_{*,t}(x;A+\delta A)-\mathcal O_{*,t}(x;A)|}{\|\delta A\|_{L^\infty}}.
\]
\end{definition}

\begin{theorem}[Exponential tail outside $B_R(x,t)$]\label{thm:tail_R}
Under the first-variation bound of Theorem~\ref{thm:first_var} and the corresponding bound for $G$ (Corollary~\ref{cor:first_var_G}),
the following exponential tail estimate holds.
There exist constants $C_{\mathrm{tail}},c_{\mathrm{tail}}>0$ such that for $*\in\{S,Q\}$,
\[
\mathrm{Lip}^{(R)}_x(\mathcal O_{*,t})\ \le\ C_{\mathrm{tail}}\,t^{-5}\,\exp(-c_{\mathrm{tail}}R^2).
\]
Consequently, if two initial data sets agree on $B_R(x,t)$, then
\[
|\mathcal O_{*,t}(x;A)-\mathcal O_{*,t}(x;A')|
\le C_{\mathrm{tail}}\,t^{-5}\,\exp(-c_{\mathrm{tail}}R^2)\,\|A-A'\|_{L^\infty}.
\]
\end{theorem}

\begin{proof}
Integrate the first functional derivative along the straight line path $A_\theta=A+\theta(A'-A)$ and use the Gaussian tail:
\[
\mathcal O_{*,t}(A')-\mathcal O_{*,t}(A)=\int_0^1 \int \frac{\delta \mathcal O_{*,t}}{\delta A(y)}(A_\theta)\,(A'-A)(y)\,dy\,d\theta.
\]
Restrict $y\in B_R^c$ and bound the derivative kernel by \eqref{eq:der_OS} or its $r=1$ version, then integrate the Gaussian tail:
$\int_{|x-y|>R\sqrt t} t^{-2}e^{-c|x-y|^2/t}dy\lesssim e^{-c'R^2}$.
\end{proof}

\begin{corollary}[Choosing $R(t)$ to match a target power $t^\eta$]\label{cor:R_choice}
Fix $\eta>0$. Choose
\[
R(t):=\sqrt{\frac{(\eta+6)\log(1/t)}{c_{\mathrm{tail}}}}\qquad (t\in(0,e^{-1}]).
\]
Then $t^{-5}e^{-c_{\mathrm{tail}}R(t)^2}\le t^{\eta+1}$.
Hence the truncation error outside $B_{R(t)}(x,t)$ is $\le C_{\mathrm{tail}} t^{\eta+1}\|A-A'\|_\infty$.
\end{corollary}

\paragraph{Interpretation for A1.}
Corollary~\ref{cor:R_choice} shows that, for any fixed target exponent $\eta$ (e.g.\ $\eta=1$ suffices for the SFTE remainder bookkeeping in Theorem~\ref{thm:ins_decomp}), one may treat $\mathcal O_{S,t}(x)$ and $\mathcal O_{Q,t}(x)$ as \emph{effectively supported} on a ball of radius $R(t)\sqrt t$,
with an error that is smaller than the desired SFTE remainder scale.
This is exactly what is needed to feed into the insertion RG map at scale $L_{j_t}\simeq \sqrt t$.

We now state the flow-localization theorem used in this submission.

\begin{theorem}[Flow quasi-locality and derivative bounds]\label{thm:flow_local_v91}
Fix $t_\ast>0$ and a finite derivative order $r_{\max}$.
Under the bounded-regime control of Proposition~\ref{ass:bounded_regime}
(verified on the good-flow chart regime by Lemma~\ref{lem:chart_persist} and absorbed in expectations by Corollary~\ref{cor:superpoly_flow})
and the Gaussian kernel bounds~\ref{ass:gaussian_kernel}
(proved in Appendix~\ref{app:kernel_parametrix} and Appendix~\ref{app:discrete_parametrix}),
the following quasi-local truncation estimate holds for $*\in\{S,Q\}$:
\begin{enumerate}[label=(L\arabic*)]
\item (\textbf{Quasi-locality}) The functional derivatives of $\mathcal O_{*,t}(x)$ with respect to the initial field at points $y$ obey Gaussian off-diagonal bounds of the form
\eqref{eq:der_OS}--\eqref{eq:der_OQ} for all $1\le r\le r_{\max}$.
\item (\textbf{Effective support}) For any $\eta>0$ there exists a radius function $R(t)\asymp \sqrt{\log(1/t)}$ such that
modifying initial data outside $B_{R(t)}(x,t)$ changes $\mathcal O_{*,t}(x)$ by at most $O(t^{\eta})$ in sup norm.
\item (\textbf{Uniformity}) All constants can be chosen independent of lattice spacing $a$ (for $t\ge c a^2$) provided the chart/boundedness hypothesis is uniform in $a$.
\end{enumerate}
\end{theorem}

\paragraph{Status in this submission.}
The bounded-regime/chart hypothesis~\ref{ass:bounded_regime} is verified on the good-flow chart event by
Lemma~\ref{lem:chart_persist} (Section~\ref{sec:good_flow_chart}), and its complement is absorbed in expectations by
Corollary~\ref{cor:superpoly_flow}.
The Gaussian kernel bounds~\ref{ass:gaussian_kernel} are proved in Appendix~\ref{app:kernel_parametrix} (continuum parametrix)
and Appendix~\ref{app:discrete_parametrix} (uniform discrete analogue).
With these ingredients in place, all subsequent uses of ``flow localization'' in the insertion RG are bookkeeping.
